% ============================================================
%  Section 6: Conclusion
% ============================================================

\section{Conclusion}

\begin{frame}{Future Work}
  \begin{columns}[T, onlytextwidth]
    \begin{column}{0.48\textwidth}
      \textbf{Model-side}
      \begin{itemize}
        \item Smaller or specialized compressors
        \item Domain fine-tuning for code, logs, and records
        \item Speculative or blockwise prediction
      \end{itemize}
    \end{column}
    \hfill
    \begin{column}{0.48\textwidth}
      \textbf{System-side}
      \begin{itemize}
        \item Native interval-capture APIs
        \item Lower-overhead stateful scheduling
        \item CPU-only paths with AVX-512 or other vector extensions
        \item Decode-aware cost models
        \item Integration with storage-tier policies
      \end{itemize}
    \end{column}
  \end{columns}
\end{frame}


\begin{frame}{Today the Fit Is Cold Database Text}
  \begin{columns}[T,onlytextwidth]
    \begin{column}{0.46\textwidth}
      \begin{block}{Plausible now}
        \begin{itemize}
          \item archival text and compliance logs
          \item write once, read rarely
          \item storage or transfer dominates
          \item structured domains such as code
        \end{itemize}
      \end{block}
    \end{column}
    \hfill
    \begin{column}{0.46\textwidth}
      \begin{block}{Not competitive now}
        \begin{itemize}
          \item hot scans and predicates
          \item high-throughput ingestion
          \item frequent updates
          \item latency-sensitive reads
        \end{itemize}
      \end{block}
    \end{column}
  \end{columns}

  \vspace{0.65em}
  \begin{tikzpicture}
    \node[draw=mDarkTeal,very thick,rounded corners=3pt,fill=mDarkTeal!7,
      minimum width=3.6cm,minimum height=1.15cm,align=center] at (2.0,0)
      {\Large\bfseries $7.7\times$ text\\[-0.1em]\scriptsize measured ratio};
    \node[draw=mDarkTeal,very thick,rounded corners=3pt,fill=mDarkTeal!7,
      minimum width=3.6cm,minimum height=1.15cm,align=center] at (6.0,0)
      {\Large\bfseries $11.2\times$ code\\[-0.1em]\scriptsize measured ratio};
    \node[draw=codecBase,very thick,rounded corners=3pt,fill=red!6,
      minimum width=3.6cm,minimum height=1.15cm,align=center] at (10.0,0)
      {\Large\bfseries $0.0345$ MB/s\\[-0.1em]\scriptsize optimized encode};
  \end{tikzpicture}
  \vfill
  \tiny Qwen2.5-0.5B, A100; ratios vary by dataset. Throughput value is \texttt{text8} encode only.
\end{frame}


\begin{frame}{Three Lessons for High-Performance LLM Compression}
  \begin{enumerate}
    \item \textbf{Measure the operator, not the model.} Generation throughput does not predict score-capture throughput.

    \vspace{0.35em}
    \item \textbf{Expose the value the consumer needs.} A target interval is cheaper to move than a vocabulary-wide score vector.

    \vspace{0.35em}
    \item \textbf{Evaluate both directions.} Encoder-only knowledge creates optimizations the decoder cannot share.
  \end{enumerate}

  \vspace{0.7em}
  \begin{columns}[c,onlytextwidth]
    \begin{column}{0.72\textwidth}
      \Large\textbf{Follow-up: Text Compression Through the Looking Glass}\\[0.35em]
      \normalsize\href{https://utndatasystems.github.io/blog/2026/text-compression/}{utndatasystems.github.io/blog/2026/text-compression/}\\[0.65em]
    \end{column}
    \begin{column}{0.24\textwidth}
      \centering
      \qrcode[height=2.2cm]{https://utndatasystems.github.io/blog/2026/text-compression/}
    \end{column}
  \end{columns}
\end{frame}
